\documentclass[journal,10pt,twocolumn]{IEEEtran}

\usepackage[utf8]{inputenc}
\usepackage[spanish]{babel}
\usepackage{graphicx}
\usepackage{amsmath}
\usepackage{booktabs}
\usepackage{cite}
\usepackage{hyperref}
\usepackage{float}
\usepackage{lipsum} % Para texto de relleno inicial

\begin{document}

\title{Comparación de Técnicas de Mejora de Contraste en Imágenes Digitales: HE, CLAHE y BHE2PL}

\author{
    \IEEEauthorblockN{Autor(es)}
    \IEEEauthorblockA{
        Facultad\\
        Asignatura: PDI - Primer Parcial\\
        Email: correo@institucion.edu
    }
}

\maketitle

\begin{abstract}
Este documento presenta un análisis cuantitativo y cualitativo de tres algoritmos de mejora de contraste en imágenes digitales: Ecualización de Histograma tradicional (HE), Ecualización de Histograma Adaptativa Limitada por Contraste (CLAHE) y Ecualización de Bi-histograma usando dos límites de meseta (BHE2PL). Se evalúa el desempeño de las técnicas utilizando métricas estandarizadas como AMBE, PSNR, entropía y contraste, sobre un conjunto de imágenes. Los resultados demuestran diferencias estadísticamente significativas en la preservación del brillo y mejora visual, validadas a través de pruebas de hipótesis rigurosas.
\end{abstract}

\begin{IEEEkeywords}
Mejora de Imagen, Ecualización de Histograma, CLAHE, BHE2PL, AMBE, PSNR.
\end{IEEEkeywords}

\section{Introducción}
\IEEEPARstart{L}{a} mejora del contraste es una etapa central en procesamiento digital de imágenes, especialmente cuando la iluminación no es uniforme o la escena presenta bajo rango dinámico. La Ecualización de Histograma (HE) es simple y rápida, pero puede sobre-realzar regiones y alterar el brillo medio. Para mitigar esas limitaciones se emplean enfoques locales, como CLAHE, y enfoques de preservación de brillo global, como BHE2PL~\cite{aquino2017bi}.

En este trabajo se comparan HE, CLAHE y BHE2PL con dos objetivos concretos: i) cuantificar el compromiso entre preservación de brillo y mejora de contraste, y ii) analizar el costo computacional de cada técnica. Para cumplirlos se combinan métricas objetivas, pruebas estadísticas pareadas y análisis visual sobre una imagen de referencia.

La estructura del artículo es la siguiente. En la Sección II se resume el marco teórico. En la Sección III se detallan el flujo experimental, las ecuaciones y los parámetros para reproducibilidad. En la Sección IV se presentan y discuten resultados cuantitativos y cualitativos. Finalmente, la Sección V expone las conclusiones.

\section{Marco Teórico}
En esta sección se describen brevemente las técnicas utilizadas.

\subsection{Ecualización de Histograma (HE)}
HE busca redistribuir las intensidades de los píxeles para lograr un histograma plano, incrementando el rango dinámico global de la imagen.

\subsection{CLAHE}
El método Contrast Limited Adaptive Histogram Equalization divide la imagen en bloques (tiles) y aplica ecualización limitando el contraste (recortando picos del histograma) para evitar la amplificación de ruido en áreas homogéneas.

\subsection{Ecualización de Bi-histograma (BHE2PL)}
Esta técnica propuesta divide el histograma en dos secciones basadas en el nivel de intensidad promedio, calculando límites de meseta para modificar de forma adaptativa cada sub-histograma y ecualizarlos independientemente, con el fin de preservar el brillo~\cite{aquino2017bi}.

\section{Metodología}
\subsection{Base de Datos y Herramientas}
Las pruebas se realizaron sobre 300 imágenes en escala de grises del conjunto Berkeley~\cite{MartinFTM01}. Se utilizaron Python 3, OpenCV, NumPy, SciPy, scikit-image y Matplotlib.

Los parámetros utilizados fueron: CLAHE con \texttt{clipLimit = 2.0} y \texttt{tileGridSize = (8,8)}; HE con la implementación estándar de OpenCV; y BHE2PL con la implementación de referencia desarrollada para este trabajo a partir de~\cite{aquino2017bi}. También se reportó tiempo de procesamiento por imagen y por método.

\subsection{Métricas de Evaluación}
Se emplearon las siguientes métricas:
\begin{itemize}
    \item \textbf{AMBE (Absolute Mean Brightness Error):} Mide la desviación del brillo promedio.
    \item \textbf{PSNR (Peak Signal-to-Noise Ratio):} Evalúa el nivel de distorsión introducido.
    \item \textbf{Contraste:} Desviación estándar de los niveles de gris.
    \item \textbf{Entropía:} Medida de la información contenida en la imagen.
    \item \textbf{SSIM (Structural Similarity Index):} Mide similitud estructural respecto a la imagen original.
    \item \textbf{Tiempo (ms):} Tiempo de ejecución por imagen para cada técnica.
\end{itemize}

Las métricas principales se calcularon mediante:
\begin{equation}
\mathrm{AMBE} = \left|\mu(I)-\mu(J)\right|
\end{equation}
\begin{equation}
\mathrm{PSNR} = 10\log_{10}\left(\frac{L^2}{\mathrm{MSE}}\right),\quad
\mathrm{MSE}=\frac{1}{MN}\sum_{x,y}\left(I(x,y)-J(x,y)\right)^2
\end{equation}
\begin{equation}
\sigma(J)=\sqrt{\frac{1}{MN}\sum_{x,y}\left(J(x,y)-\mu(J)\right)^2}
\end{equation}
\begin{equation}
H(J)=-\sum_{k=0}^{L-1} p_k\log_2 p_k
\end{equation}
donde $I$ es la imagen original, $J$ la imagen mejorada, $\mu$ la media de intensidad, $\sigma$ el contraste por desviación estándar y $H$ la entropía.

\subsection{Análisis Estadístico}
Para evaluar diferencias entre métodos se aplicó la prueba $t$ de Student pareada por imagen. Se utilizó un nivel de significancia $\alpha=0.05$. El valor $p$ indica qué tan compatible es la diferencia observada con la hipótesis nula de igualdad entre métodos: cuanto menor es $p$, mayor evidencia de diferencia estadística. Para AMBE, PSNR, SSIM y tiempo se compararon HE, CLAHE y BHE2PL. Para contraste y entropía se incluyó también la imagen original.

\section{Resultados y Discusión}
\subsection{Análisis Cuantitativo}
La Tabla~\ref{tab:global_results} resume los promedios globales calculados en el script de evaluación, reportando media $\pm$ desviación estándar.

\begin{table}[H]
\caption{Promedios globales (media $\pm$ desvío estándar) para 300 imágenes}
\label{tab:global_results}
\centering
\resizebox{\columnwidth}{!}{%
\begin{tabular}{lccccc}
\toprule
Método & T (ms) & AMBE & PSNR & Entr. & Cont. \\
\midrule
Original & -- & -- & -- & 7.127429 $\pm$ 0.534419 & 50.061337 $\pm$ 15.067624 \\
HE & \textbf{0.088105 $\pm$ 0.017721} & 27.804682 $\pm$ 20.683897 & 16.117495 $\pm$ 4.109758 & 6.931329 $\pm$ 0.543245 & \textbf{73.307259 $\pm$ 2.032467} \\
CLAHE & 0.285640 $\pm$ 0.079161 & 12.020926 $\pm$ 7.570239 & 20.079830 $\pm$ 2.069813 & \textbf{7.553357 $\pm$ 0.377354} & 57.584592 $\pm$ 11.488977 \\
BHE2PL & 1.429696 $\pm$ 0.189326 & \textbf{1.315239 $\pm$ 0.934198} & \textbf{42.159020 $\pm$ 19.368343} & 7.075197 $\pm$ 0.526298 & 51.321762 $\pm$ 14.434098 \\
\bottomrule
\end{tabular}%
}
\vspace{2pt}
\footnotesize{Los mejores resultados por métrica se resaltan en negrita.}
\end{table}

La Tabla~\ref{tab:global_results} muestra que BHE2PL logra el mejor AMBE (menor error de brillo) y el mejor PSNR, mientras que HE es el método más rápido y el que más incrementa contraste global. CLAHE alcanza la mayor entropía promedio. Este comportamiento es coherente con la naturaleza de cada método: HE prioriza expansión global del histograma, CLAHE enfatiza realce local y BHE2PL balancea realce con preservación de brillo.

En términos estadísticos, las comparaciones pareadas arrojaron valores $p$ extremadamente bajos para AMBE y PSNR entre pares de métodos (por ejemplo, AMBE HE vs CLAHE: $1.7857\times10^{-43}$), lo que respalda que las diferencias observadas no se deben al azar en este conjunto de datos.

\subsection{Análisis Visual Cualitativo}
La Fig.~\ref{fig:ref_images} presenta la imagen de referencia original y los resultados de cada técnica. La Fig.~\ref{fig:ref_hist} muestra los histogramas correspondientes. Ambas figuras permiten relacionar la evidencia visual con los resultados de la Tabla~\ref{tab:global_results}.

\begin{figure}[H]
\centering
\includegraphics[width=0.48\columnwidth]{output/referencia_original.png}
\includegraphics[width=0.48\columnwidth]{output/referencia_he.png}\\
\includegraphics[width=0.48\columnwidth]{output/referencia_clahe.png}
\includegraphics[width=0.48\columnwidth]{output/referencia_bhe2pl.png}
\caption{Comparación visual para una imagen de referencia (arriba izquierda: Original, arriba derecha: HE, abajo izquierda: CLAHE, abajo derecha: BHE2PL).}
\label{fig:ref_images}
\end{figure}

\begin{figure}[H]
\centering
\includegraphics[width=\columnwidth]{output/referencia_histogramas.png}
\caption{Histogramas de la imagen de referencia y de las técnicas evaluadas.}
\label{fig:ref_hist}
\end{figure}

Visualmente, HE produce el mayor realce global y puede saturar regiones brillantes u oscuras. CLAHE mejora detalles locales, aunque en algunos casos aumenta la apariencia de ruido. BHE2PL ofrece una transición tonal más conservadora y consistente con la menor variación de brillo medida por AMBE.

\section{Conclusiones}
Los experimentos cumplieron los objetivos planteados. Primero, se confirmó cuantitativamente que BHE2PL preserva mejor el brillo que HE y CLAHE (AMBE mínimo), manteniendo además PSNR superior. Segundo, se verificó que HE es el método con menor costo temporal promedio, por lo que resulta útil cuando la prioridad es velocidad.

La discusión conjunta de Tabla~\ref{tab:global_results}, Fig.~\ref{fig:ref_images} y Fig.~\ref{fig:ref_hist} muestra que no existe un método universalmente óptimo: la selección depende del objetivo de aplicación. Si se prioriza preservación de brillo y fidelidad, BHE2PL es la opción más robusta. Si se prioriza realce local, CLAHE ofrece mejor compromiso visual. Si se prioriza costo computacional y contraste global, HE es conveniente.

\bibliographystyle{IEEEtran}
\bibliography{referencias}

\end{document}
